\documentclass[11pt,a4paper]{article}
\usepackage[utf8]{inputenc}
\usepackage[T1]{fontenc}
\usepackage[polish]{babel}
\usepackage{lmodern}
\usepackage[margin=27mm]{geometry}
\usepackage{microtype}
\usepackage{amsmath}
\usepackage[hidelinks]{hyperref}
\title{Dwuwymiarowa mapa doliny niesamowitości: realizm wyglądu i naturalność zachowania}
\author{Mateusz Pomianek}
\date{Szkic naukowy, wersja 0.4.0, 27 listopada 2026}
\begin{document}
\maketitle
\begin{abstract}
Czy realizm wyglądu postaci oraz naturalność jej ruchu niezależnie i interakcyjnie determinują oceny niesamowitości? Przedstawiono prospektywny projekt badania oparty na wspólnej platformie ST-08, wykorzystujący archiwalne nagrania zachowań odtwarzane na syntetycznych postaciach o kontrolowanym wyglądzie. Pierwszorzędowy wskaźnik: Interakcja A i B na skali latentnej modelu porządkowego; marginalna powierzchnia ocen jako odrębny estymand. Badanie nie dostarczyło dotąd danych ani wyników empirycznych.
\end{abstract}
\section{Wprowadzenie}
Zależność między wyglądem jednostki a widocznym przebiegiem jej zachowania jest zasadniczym elementem społecznej percepcji. W kontaktach z robotami społecznymi i realistycznymi awatarami nieoczekiwana mimika, gest, kontakt wzrokowy i dynamika wypowiedzi mogą wpływać na reakcje odbiorcy. Nie oznacza to jednak, że każda atypowa ekspresja musi być oceniana negatywnie, ani że jest uprawnione utożsamianie ekspresji z intensywnością doświadczanych emocji.

Projekt stanowiska ST-08 umożliwi oddzielenie zachowania rzeczywistego człowieka, wyodrębnionego z istniejących nagrań interakcji, od syntetycznej formy jego reprezentacji. Pozwala to porównywać te same zdarzenia w kontrolowanym wyglądzie awatara. Nie można jednak przyjąć, że retargeting automatycznie zachowuje wszystkie własności zachowania. Jego dokładność i sposób deformowania sygnałów muszą być sprawdzone niezależnie od właściwego eksperymentu.
\section{Podstawy teoretyczne}
Hipoteza doliny niesamowitości zakłada możliwość niemonotonicznej zależności między podobieństwem do człowieka a reakcją afektywną (Mori et al., 2012). Przegląd Kätsyri et al. (2015) wskazuje niedopasowanie percepcyjne jako jeden z możliwych mechanizmów, nie jako uniwersalne wyjaśnienie wszystkich efektów. Badania nad predykcją działania (Saygin et al., 2012) uzasadniają odrębne rozpatrywanie wyglądu i ruchu, lecz nie dowodzą, że ich interakcja musi generować lokalne maksimum subiektywnego dyskomfortu.

Projekt rozróżnia trzy konkurencyjne struktury: addytywność, interakcję bez maksimum wewnętrznego oraz interakcję z maksimum wewnętrznym. Rozbieżność między wyglądem i zachowaniem jest hipotezą mechanistyczną, której nie można utożsamiać z samym statystycznym odrzuceniem addytywności. Oceny niesamowitości należy odróżniać od atrakcyjności, znajomości i przypisywanej sprawczości. Badanie nie obejmuje bezpośredniego pomiaru neuronalnego błędu predykcji.

\section{Cel i hipotezy}
Czy realizm wyglądu postaci oraz naturalność jej ruchu niezależnie i interakcyjnie determinują oceny niesamowitości?
\paragraph{H1.} Realizm wyglądu i naturalność ruchu oddziałują na subiektywną ocenę niesamowitości ponad ich niezależne efekty.
\paragraph{H2.} Powierzchnia ocen może mieć lokalne maksima dyskomfortu, lecz istnienie takiej struktury nie jest przyjmowane z góry.
Przed badaniem potwierdzającym konieczne jest ustalenie estymandu, sposobu kodowania efektu, kryteriów jakości, mocy oraz procedury radzenia sobie z wielokrotnymi porównaniami.
\section{Materiał i metody}
Wspólne zwalidowane klipy ST-08 z niezależną manipulacją wyglądu A i kinematyki B. Zależności czasowe pierwotnego zachowania, kadr i kontekst są kontrolowane.

Randomizowany eksperyment ekranowy z kilkoma poziomami A i B; kalibracja percepcyjna, kilka postaci i źródłowych interakcji, efekty wewnątrzosobnicze.

Pierwszorzędowy wskaźnik: Interakcja A i B na skali latentnej modelu porządkowego; marginalna powierzchnia ocen jako odrębny estymand. Jednostką analizy jest ocena konkretnej sekwencji przez konkretnego obserwatora; należy uwzględnić zależność prób od źródłowej osoby, scenariusza i postaci. Zbieranie odpowiedzi behawioralnych nie wymaga EEG ani okulografii. Szacowana orientacja spojrzenia w archiwalnym filmie nie jest laboratoryjnym pomiarem eye trackingu.
\section{Plan analizy}
Podstawowy model konfirmacyjny jest kumulacyjnym modelem porządkowym z interakcją czynników A i B. Modele GAMM pozostają eksploracyjne.

Zamrożenie estymandu i planu porównań nastąpi po pilotażu technicznym i przed dostępem do danych konfirmacyjnych. Należy zachować oddzielnie analizy eksploracyjne, raportować przedziały niepewności, wielkości efektów oraz reguły wyłączenia bodźców z błędną animacją. Trafność konstruktu wymaga niezależnej kalibracji realizmu wyglądu i percepcyjnej naturalności zachowania.
\section{Operacjonalizacja i projekt eksperymentu}
Proponuje się układ czynnikowy 4 na 4, z czterema poziomami realizmu wyglądu i czterema poziomami naturalności ruchu. Poziomy są robocze i wymagają niezależnej kalibracji. Manipulacja ruchem powinna dotyczyć parametrów animacji, takich jak ciągłość i koordynacja czasowa, a nie przypisywanej osobie diagnozy. Treść sceny, dźwięk, czas ekspozycji i kadr pozostają stałe. Każda komórka powinna obejmować wiele niezależnych źródeł zachowania i tożsamości awatara. Zrównoważone bloki niepełne ograniczą zmęczenie oraz wielokrotną ekspozycję na ten sam materiał źródłowy.

Pilotaż z niezależnymi obserwatorami oszacuje rozdzielność manipulacji, jakość retargetingu i zrozumiałość skali niesamowitości. Kryteria akceptacji bodźców oraz minimalna istotna różnica zostaną ustalone przed badaniem głównym. Nie należy wykluczać bodźców na podstawie ich ocen niesamowitości w próbie konfirmacyjnej.

\section{Estymandy i analiza konfirmacyjna}
Niech $\mu_{ij}$ oznacza marginalną oczekiwaną ocenę niesamowitości dla poziomów $i,j$, uśrednioną względem docelowej populacji obserwatorów i bodźców. Interakcję definiuje się przez
\[
D_{ij}=\mu_{ij}-\mu_{i\cdot}-\mu_{\cdot j}+\mu_{\cdot\cdot}.
\]
Hipoteza zerowa zakłada $D_{ij}=0$ dla wszystkich komórek. Podstawowy test obejmuje dziewięć stopni swobody interakcji w pełnym układzie 4 na 4. Model mieszany uwzględni skrzyżowane efekty obserwatora, źródła zachowania i tożsamości awatara, o ile projekt umożliwi ich identyfikację. Dla odpowiedzi porządkowej zastosowany zostanie model kumulacyjny; kontrasty marginalne na skali oczekiwanej oceny będą odróżniane od interakcji na skali latentnej. Szczegóły testu zostaną zamrożone po pilotażu, bez dostępu do danych głównych.

Hipoteza lokalnego maksimum wymaga porównania komórki wewnętrznej z jej czterema bezpośrednimi sąsiadami. Wnioskowanie o maksimum wymaga dodatnich dolnych granic jednoczesnych przedziałów dla wszystkich czterech kontrastów. Korekta obejmie wszystkie rozpatrywane komórki wewnętrzne. Interpolowana powierzchnia GAMM pozostaje analizą eksploracyjną, dopóki kalibracja nie uzasadni ciągłej skali obu wymiarów.

\section{Plan symulacji i kryteria wykonalności}
Symulacja Monte Carlo powinna uwzględniać scenariusz zerowy, efekty addytywne, interakcję bez lokalnego maksimum oraz lokalne maksimum. Zmieniane będą liczba obserwatorów, liczba niezależnych bodźców, wariancja efektów losowych, rozkład odpowiedzi i odsetek braków. Raport obejmie błąd pierwszego rodzaju, moc, pokrycie przedziałów, zbieżność oraz niepewność Monte Carlo. Liczebność próby nie jest obecnie ustalona. Warunkiem rozpoczęcia rekrutacji są prawa do materiałów, kwalifikacja etyczna, walidacja bodźców i zamrożenie protokołu.

\section{Wyniki}
Nie zebrano jeszcze danych empirycznych i nie wykonano analiz potwierdzających. W niniejszym opracowaniu nie przypisuje się wartości statystycznych, wyników skal, skuteczności diagnostycznej ani dowodów występowania doliny niesamowitości.
\section{Dyskusja i ograniczenia}
Jest to mapa podstawowego efektu, nie test przyczynowego wpływu oczekiwań z P03 ani naturalnych profili komunikacji z P04. Prawa do korzystania z oryginalnych nagrań i rozpowszechniania pochodnych animacji nie wynikają automatycznie z publicznej dostępności bazy. Rekonstrukcja mimiki i ruchu z pojedynczej kamery może zawierać systematyczne błędy. Dla badań z udziałem nowych obserwatorów konieczne są właściwa kwalifikacja etyczna, ochrona danych i jasno określona procedura zgody.

Wyniki negatywne oraz niejednoznaczne mają znaczenie naukowe, jeśli kontrolowano jakość bodźców, precyzję analizy i heterogeniczność kontekstu. Szczególnie w przypadku różnorodności komunikacji nie można wywodzić wewnętrznych stanów psychicznych z nietypowej ekspresji widocznej na twarzy ani sprowadzać zróżnicowania do jednolitej cechy klinicznej.
\section{Wnioski}
Zaproponowano odrębny problem naukowy i sposób jego przyszłego falsyfikowania. Uprawnienia do nagrań, walidacja rekonstrukcji zachowania, prerejestracja i badania obserwatorów pozostają do wykonania. Powstanie szkicu LaTeX nie stanowi dowodu gotowości publikacyjnej.
\section*{Literatura}
\begin{description}
\item Kätsyri, J., Förger, K., Mäkäräinen, M., \& Takala, T. (2015). A review of empirical evidence on different uncanny valley hypotheses: Support for perceptual mismatch as one road to the valley of eeriness. \textit{Frontiers in Psychology, 6}, Article 390. \url{https://doi.org/10.3389/fpsyg.2015.00390}
\item Mori, M., MacDorman, K. F., \& Kageki, N. (2012). The uncanny valley [From the field]. \textit{IEEE Robotics \& Automation Magazine, 19}(2), 98--100. \url{https://doi.org/10.1109/MRA.2012.2192811}
\item Saygin, A. P., Chaminade, T., Ishiguro, H., Driver, J., \& Frith, C. (2012). The thing that should not be: Predictive coding and the uncanny valley in perceiving human and humanoid robot actions. \textit{Social Cognitive and Affective Neuroscience, 7}(4), 413--422. \url{https://doi.org/10.1093/scan/nsr025}
\end{description}
\end{document}
